\subsection{Radial convolution after two inverse lifts}
\label{subsec:two-lift-convolution}

The preceding pressure cancellations reduce the remaining estimates to products of analytic multipliers with regularized derivatives. A coefficientwise bound on an individual Taylor monomial is not enough for this purpose: the leading profiles have infinitely many radial coefficients, and differentiation distributes among them. We establish here an operator estimate that includes the complete radial convolution, without imposing a finite-degree cutoff. It also records which normal-form reduction is still needed to apply the estimate to the actual pressure source.

\begin{theorem}[A uniform multiplication constant]\label{thm:algebra64}
For the doubled coefficient weight \eqref{eq:weight}, the induced norm \eqref{eq:norm} satisfies
\begin{equation}\label{eq:algebra64}
 \|fg\|_{R,\rho;E}\le 64\|f\|_{R,\rho;E}\|g\|_{R,\rho;E}
\end{equation}
for all coefficient series for which these norms are finite. Moreover, $I_mI_p\subseteq I_{m+p}$, with the same constant. This improves the admissible coefficient-algebra constant in Appendix B.1 of \cite{newton2026}.
\end{theorem}
\begin{proof}
Write $N=i+j$ and $b=u+v$. Leibniz's rule and Vandermonde's inequality give
\[
 \binom b u\frac{w_{i,u}w_{j,v}}{w_{N,b}}
 \le \frac{(N+1)^2}{(i+1)^2(j+1)^2}
       \frac{(b+1)^2}{(u+1)^2(v+1)^2},
\]
since $\binom{2i+u}{u}\binom{2j+v}{v}\le\binom{2N+b}{b}$.
For each $N\ge0$,
\[
 (N+1)^2\sum_{i=0}^N\frac1{(i+1)^2(N-i+1)^2}
 \le4\sum_{r=1}^{N+1}\frac1{r^2}\le8.
\]
Indeed $(N+1)^2\le 2[(i+1)^2+(N-i+1)^2]$, while
$\sum_{r=1}^{N+1}r^{-2}\le1+\sum_{r=2}^{N+1}[r(r-1)]^{-1}<2$.
The same estimate applies to the Leibniz derivative index $b$. Their product is $64$, uniformly in both indices. The valuation assertion follows directly from coefficient convolution.
\end{proof}

The next result allows multiplication by a nonpolynomial analytic profile \emph{before} the radial inverses act. Let $D_X=X\partial_X$ and $n\ge1$. For $0\le j\le3$ and $F\in I_n$, define the two-lift jet operator
\begin{equation}\label{eq:regularized-jet}
 (\mathcal R_{n,j}F)_{a+2}
 =\frac{(a+n+2)^{3-j}}{(a+1)^2(a+2)^2}\,\partial_\eta^j F_a,
 \qquad a\ge n.
\end{equation}
For a multiplier $g=\sum_{c\ge0}g_c(\eta)X^c$, define the convolution-inverse operator coefficientwise by
\begin{equation}\label{eq:radial-convolution}
 (\mathcal C_{n,j}(g,F))_{N}
 =\sum_{\substack{a\ge n,\ c\ge0\\a+c+2=N}}
 \frac{(N+n)^{3-j}}{(N-1)^2N^2}\,
 g_c\,\partial_\eta^jF_a,
 \qquad N\ge n+2.
\end{equation}
No restriction is placed on the number of nonzero Taylor coefficients of $g$.

\begin{theorem}[Tame analytic radial-convolution smoothing]\label{thm:convolution-smoothing}
Let $0\le j\le3$, $n\ge1$, $R,\rho>0$, and put
\begin{equation}\label{eq:tame-constant}
 D_j(R,\rho)=\frac{675\,2^{4-j}R^2}{(j+1)^2\rho^j}.
\end{equation}
Then
\begin{equation}\label{eq:tame-basic}
 \|\mathcal R_{n,j}F\|_{R,\rho;E}
 \le\frac{D_j(R,\rho)}{n+1}\|F\|_{R,\rho;E}.
\end{equation}
If $g,D_Xg,\ldots,D_X^{3-j}g$ belong to the same weighted space, then $\mathcal C_{n,j}(g,F)\in I_{n+2}$ and
\begin{equation}\label{eq:tame-convolution}
 \|\mathcal C_{n,j}(g,F)\|_{R,\rho;E}
 \le\frac{64D_j(R,\rho)}{n+1}
 \left(\sum_{s=0}^{3-j}\binom{3-j}{s}\|D_X^sg\|_{R,\rho;E}\right)
 \|F\|_{R,\rho;E}.
\end{equation}
Both constants are independent of $n$ and of the radial Taylor support of $g$.
\end{theorem}
\begin{proof}
Set $t=a+n+2$. Since $a\ge n$, one has
$2(n+1)\le t\le2(a+1)$. The exact fourth-jet weight identity
\eqref{eq:fourjet-exact} implies, for $0\le j\le3$,
\[
 \frac{t^{4-j}w_{a,b+j}}{(a+1)^2(a+2)^2w_{a+2,b}}
 \le \frac{54R^2}{\rho^j}\frac{25\,2^{4-j}}{(j+1)^2}.
\]
To obtain this, divide the $j$-jet ratio by the fourth-jet ratio. The quotient is
$\rho^{4-j}(b+5)^2/(b+j+1)^2$ times the reciprocal of
$\prod_{m=j+1}^4(2a+b+m)$, where the first fraction is at most
$25\rho^{4-j}/(j+1)^2$ and each denominator is at least $a+1$.
Dividing by $t\ge2(n+1)$ gives \eqref{eq:tame-basic}, with
$54\cdot25/2=675$.

For the full convolution, write $p=3-j$, $N=a+c+2$, and
$t=a+n+2$. The elementary facts
\[
 (N+n)^p=(t+c)^p=\sum_{s=0}^p\binom p s t^{p-s}c^s,
 \qquad (N-1)^2N^2\ge(a+1)^2(a+2)^2,
\]
show that each coefficient and each parameter derivative of
$\mathcal C_{n,j}(g,F)$ is dominated term by term by the nonnegative
Leibniz majorants of
\[\sum_{s=0}^p\binom p s(D_X^sg)\,\mathcal R_{n,j}F.\]
Here we used $t\ge1$ to replace $t^{p-s}$ by $t^p$; no derivative is moved between $g$ and $F$ without the corresponding Leibniz terms. Apply the multiplication estimate
\eqref{eq:algebra64} to these majorants and then
\eqref{eq:tame-basic}. This proves \eqref{eq:tame-convolution}.
\end{proof}

\paragraph{The highest leading-profile pressure jet.}
The source-native formula \eqref{eq:leading-pressure-operator} allows the highest input derivative to be identified without freezing any leading-profile coefficient. Fix $k\ge1$, $a\ge k$ and $c\ge0$, set $r=a+1$, $m=a+c+1$, $N=m+1$, and let $u_c$ be the $X^c$ coefficient of $U_0$. For one input coefficient $F_a$, the radial-velocity identity \eqref{eq:leading-W} reads
\[
 W_{k,a}=w_0F_a+w_1\partial_\eta F_a,\quad
 w_0=B\frac{r-D-\lambda_k}{r},\qquad w_1=-\frac Gr.
\]
Only the $U_0Z_{0,k}V_k$ contribution can produce three derivatives after the final pressure reconstruction. The full source operator's third-jet coefficient is exactly
\begin{equation}\label{eq:highest-leading-jet}
 \bigl(\mathcal P_{k+1}^{\mathrm{lead}}F\bigr)_{a+c+2}\big|_{\partial_\eta^3F_a}
 =\frac{u_cG^3}{4r m N^2}\,\partial_\eta^3F_a.
\end{equation}
This is an identity of variable-coefficient operators, including the radial degree $c$ of the leading profile.

\begin{corollary}[Uniform leading-profile highest-jet channel]\label{cor:leading-highest}
If $g=\tfrac14G^3(U_0+D_XU_0)$ belongs to the weighted algebra, the complete third-derivative component of \eqref{eq:leading-pressure-operator}, summed over all $a,c$, satisfies
\[
 \|\mathcal P_{k+1}^{\mathrm{lead},(3)}F\|_{R,\rho;E}
 \le\frac{64D_3(R,\rho)}{k+1}\|g\|_{R,\rho;E}\|F\|_{R,\rho;E},\qquad F\in I_k.
\]
\end{corollary}
\begin{proof}
The coefficient of $\partial_\eta^2F_a$ in the integrand's $U_0Z_{0,k}V_k$ term is $u_cGw_1=-u_cG^2/r$. The radial primitive supplies $1/m$, and the outer $G\partial_\eta$ followed by $-J_1/4$ yields \eqref{eq:highest-leading-jet}. In the model operator \eqref{eq:radial-convolution} with $j=3$ the denominator is $m^2N^2$. Since $m/r=1+c/r\le1+c$,
\[
 \frac{|u_cG^3|}{4rmN^2}
 \le\frac{(1+c)|u_cG^3|}{4m^2N^2}.
\]
The multiplier on the right is precisely the $X^c$ coefficient of $g$. The same domination holds termwise after any $\eta$-derivatives, because $m/r$ is independent of $\eta$. Apply Theorem~\ref{thm:convolution-smoothing} with $j=3$.
\end{proof}

This result is directly applicable whenever a term of the source operator can be normal-ordered into a finite sum of the coefficients in \eqref{eq:radial-convolution}, with multipliers in the stated common analytic space. In particular, it removes the need to truncate a leading profile's radial series and controls the additional factors of $c$ generated by $D_X$ commutators. The specific highest input-jet channel in the leading-profile pressure source is covered by Corollary~\ref{cor:leading-highest}. The next subsection closes the three remaining leading-profile input-jet channels. This abstract theorem alone does \emph{not} prove that the complete pressure operators \eqref{eq:leading-pressure-operator} and \eqref{eq:quad-reconstruct} have yet been reduced to such finite sums with a uniform multiplier budget. Their exact symbolic normal forms include intermediate radial averages and, in the quadratic case, two separate unknown inputs. The required source-specific reduction and the bilinear estimate for $H_k$ remain the next proof obligations. The complete recurrence (53) and Theorem~1.1 therefore remain open.
